\documentclass[11pt,letterpaper]{article}

% ======= PACKAGES =======
\usepackage[utf8]{inputenc}
\usepackage{geometry}
\geometry{margin=1in}
\usepackage{setspace}
\setstretch{1.5}
\usepackage{fancyhdr}
\usepackage{xcolor}
\usepackage{titlesec}
\usepackage{lastpage}
\usepackage{hyperref}
\usepackage{mathpazo}
\usepackage{graphicx}
\usepackage{subcaption}
\linespread{1.5}

% ======= HEADER & FOOTER =======
\pagestyle{fancy}
\fancyhf{}

% Header (faded + smaller)
\fancyhead[L]{\footnotesize\textcolor{gray!90}{\textit{Biographical Sketch}}}
\fancyhead[R]{\footnotesize\textcolor{gray!90}{\textit{Mostafa Rezaali}}}

% Footer (faded)
\fancyfoot[L]{\footnotesize\textcolor{gray!90}{\textit{NSF AGS-PRF -- University of Florida}}}
\fancyfoot[R]{\footnotesize\textcolor{gray!90}{Page \thepage}}

\renewcommand{\headrulewidth}{0.4pt}
\renewcommand{\footrulewidth}{0.4pt}

% ======= PARAGRAPH STYLE =======
\setlength{\parindent}{15pt}
\setlength{\parskip}{6pt}
\titlespacing*{\section}{0pt}{0.55em}{0.2em}
\titlespacing*{\subsection}{0pt}{0.45em}{0.15em}

% ======= DOCUMENT BEGINS =======
\begin{document}

% ======= TITLE BLOCK =======
\begin{center}
    \vspace*{-1.5cm}
    {\Large\bfseries Postdoctoral Fellowship Application} \\[0.2cm]
    {\small\textit{NSF Atmospheric and Geospace Sciences Postdoctoral Research Fellowship}} \\[0.05cm]
    {\small\textit{Department of Geography, University of Florida}} \\[0.1cm]
    \rule{0.6\textwidth}{0.4pt}
\end{center}
\vspace{-0.5cm}

\noindent
Mostafa Rezaali \\
Ph.D. Candidate, Climate Science (Geography) \\
University of Florida \\
\href{mailto:mostafarezaali@gmail.com}{mostafarezaali@gmail.com} \quad $\mid$ \quad (352) 709-5350 \\[0.3cm]
{\large\bfseries Biographical Sketch}\\[-0.15cm]
\rule{\textwidth}{0.4pt}
\section*{Identifying Information}
Mostafa Rezaali is a Ph.D. Candidate in Climate Science (Geography) at the University of Florida and a lawful permanent resident of the United States. His research develops artificial-intelligence methods for climate extremes and environmental prediction, with emphasis on heat waves, flash droughts, spatiotemporal climate forecasting, air-quality modeling, and uncertainty-aware decision support. His NSF AGS-PRF work is centered on conditional flow matching (CFM) and graph-based atmospheric representations for probabilistic prediction of subseasonal U.S. heat extremes.

Rezaali's training spans climate science, geography, civil and environmental engineering, atmospheric sciences, machine learning, and geospatial data science. This interdisciplinary preparation positions him to translate recent advances in graph neural networks, generative modeling, and probabilistic forecast verification into tools for atmospheric science, climate-risk assessment, and public-interest forecasting.

\section*{Professional Preparation}
\begin{tabular}{p{0.30\textwidth}p{0.27\textwidth}p{0.18\textwidth}p{0.15\textwidth}}
\textbf{Institution} & \textbf{Location} & \textbf{Degree/Area} & \textbf{Dates}\\
University of Florida & Gainesville, FL & Ph.D., Climate Science/Geography & 2022--present\\
University of Florida & Gainesville, FL & Graduate Certificate, Atmospheric Sciences & 2022--present\\
Qom University of Technology & Qom, Iran & M.Sc., Civil and Environmental Engineering & 2016--2018\\
IAUKHSH & Isfahan, Iran & B.Sc., Civil and Environmental Engineering & 2011--2016\\
\end{tabular}

\section*{Appointments and Positions}
\begin{itemize}
\item Graduate Research Assistant, Department of Geography, University of Florida, 2022--present.
\item NSF LEAP REU Mentor, 2025.
\item Invited Lecturer, Deep Learning Applications in Climate Science, University of Florida, 2025.
\item Invited Lecturer, Climate Change Impacts on Future Plant Distributions, University of Florida, 2025.
\item M.Sc. Student Advisor, Alborz University of Medical Sciences, 2021--present.
\item Research and modeling contributor to projects involving machine learning, climate extremes, water-demand forecasting, air-quality prediction, and environmental-health risk.
\end{itemize}

\section*{Research Expertise and Technical Preparation}
Rezaali's research portfolio emphasizes data-driven modeling of high-dimensional environmental systems. He has worked with large gridded climate and atmospheric datasets, including daily temperature fields, atmospheric circulation variables, surface and land-state predictors, and environmental-monitoring data. His computational experience includes Python, PyTorch-style deep learning workflows, R, MATLAB, ArcGIS/ArcPy, Linux environments, \LaTeX, NetCDF/GRIB processing, spatial statistics, model evaluation, and visualization.

For the proposed fellowship, this technical preparation is directly relevant to the development of probabilistic graph models on an icosahedral mesh. The planned CFM framework builds on Rezaali's prior experience with ensemble deep learning, residual prediction, cross-validation over historical climate records, and forecast verification against persistence and climatology baselines. His prior climate and environmental applications provide a strong foundation for evaluating whether modern generative learning systems can extract physically meaningful, calibrated predictability from atmospheric precursor fields.

\section*{Products Most Closely Related to the Proposed Project}
\begin{enumerate}
\item Rezaali, M., Jahangir, M. S., Fouladi-Fard, R., and Keellings, D. (2024). An ensemble deep learning approach to spatiotemporal tropospheric ozone forecasting: A case study of Tehran, Iran. \textit{Urban Climate}, 55, 101950.
\item Narayanan, A., Rezaali, M., Bunting, E. L., and Keellings, D. (2025). It's getting hot in here: Spatial impact of humidity on heat wave severity in the US. \textit{Science of The Total Environment}, 963, 178397.
\item Rezaali, M., Quilty, J., and Karimi, A. (2021). Probabilistic urban water demand forecasting using wavelet-based machine learning models. \textit{Journal of Hydrology}, 600, 126358.
\item Rezaali, M., Fouladi-Fard, R., Mojarad, H., Sorooshian, A., Mahdinia, M., et al. (2021). A wavelet-based random forest approach for indoor BTEX spatiotemporal modeling and health risk assessment. \textit{Environmental Science and Pollution Research}, 28, 22522--22535.
\item Rezaali, M., Fouladi-Fard, R., and Karimi, A. (2023). Performance of TANN, NARX, and GMDHT models for urban water demand forecasting: A case study in a residential complex in Qom, Iran. \textit{Avicenna Journal of Environmental Health Engineering}, 10(2), 85--97.
\end{enumerate}

\section*{Other Significant Products}
\begin{enumerate}
\item Rezaali, M., Fouladi-Fard, R., O'Shaughnessy, P., Naddafi, K., and Karimi, A. (2025). Assessment of AERMOD and ADMS for NOx dispersion modeling with a combination of line and point sources. \textit{Stochastic Environmental Research and Risk Assessment}, 1--15.
\item Rezaali, M., Fouladi-Fard, R., and Karimi, A. (2020). Identification of temporal and spatial patterns of river water quality parameters using NLPCA and multivariate statistical techniques. \textit{International Journal of Environmental Science and Technology}, 17, 2977--2994.
\item Rezaali, M., and Fouladi-Fard, R. (2021). Aerosolized SARS-CoV-2 exposure assessment: dispersion modeling with AERMOD. \textit{Journal of Environmental Health Science and Engineering}, 19, 285--293.
\end{enumerate}

\section*{Graduate Training, Mentoring, and Educational Contributions}
Rezaali has integrated research with mentoring and teaching. As an NSF LEAP REU mentor, he helped undergraduate researchers engage with climate and environmental data science, including problem framing, data preparation, model development, interpretation, and communication of results. His guest lectures at the University of Florida introduced students to deep learning applications in climate science and to modeling climate-change impacts on plant distributions, linking technical methods to applied climate questions.

His mentoring philosophy emphasizes practical reproducibility, transparent uncertainty communication, and the responsible use of AI in environmental decision-making. These activities support the broader educational goals of the AGS-PRF by preparing him to lead interdisciplinary research groups that connect atmospheric science, geospatial data science, and public-interest climate services.

\section*{Synergistic and Service Highlights}
Reviewer service includes \textit{Scientific Reports}, \textit{Journal of Hydrology}, and \textit{Springer Nature Applied Sciences}. Teaching and mentoring activities include graduate guest lectures in climate-data science and mentoring of undergraduate and master's students on neural-network applications to climate and environmental hazards. Rezaali's broader professional contributions are unified by a commitment to making AI-based climate prediction more interpretable, reproducible, and useful for communities exposed to extreme heat, drought, and environmental-health risks.

\end{document}
